\documentclass[11pt]{article}
\usepackage{ap-psych}
\title{Unit 5: Mental and Physical Health}
\author{Paul Wang}
\date{November 2026}
\docpart{Unit 5 \textperiodcentered\ Study Guide \textperiodcentered\ 15--25\% of MCQ}

\begin{document}
\maketitle

\begin{abstract}
Unit 5 applies the whole course to health. It covers stress and coping (5.1), positive psychology (5.2), how psychological disorders are defined, explained, and classified (5.3), the major categories of disorders in the DSM (5.4), and psychological and biomedical treatments (5.5). Exam items usually describe a person's symptoms and ask for the disorder, or describe a therapist's technique and ask for the approach, so the guide lists the distinguishing symptom or technique for each. The descriptions here are for exam study, not diagnosis. Each section ends with check questions; the key is on the last page.
\end{abstract}

\scoringline{Check questions: \blank\ / \qtotal \qquad Target: 20+ \qquad Disorders I confuse: \blank[10em]}

\section{Health psychology and stress (5.1)}

\textbf{Health psychology} studies how psychological, behavioral, and social factors affect physical health (the \textbf{biopsychosocial} approach). \textbf{Stress} is the process of appraising and responding to events seen as threatening or challenging. \textbf{Eustress} is positive, motivating stress; \textbf{distress} is harmful stress.

\textbf{Stressors}: \textbf{catastrophes} (disasters, war), \textbf{significant life changes} (a death, a move, divorce), \textbf{daily hassles} (traffic, deadlines), and \textbf{adverse childhood experiences} (ACEs).

\begin{compare}{Stage}{Selye's General Adaptation Syndrome}{Body}
1. Alarm & Sympathetic nervous system activates; resources mobilized. & Stress hormones (adrenaline, cortisol) surge. \\
2. Resistance & Body copes with the stressor at high arousal. & Temperature, blood pressure, respiration stay elevated. \\
3. Exhaustion & Reserves are depleted under prolonged stress. & Greater vulnerability to illness. \\
\end{compare}
\textbf{Tend-and-befriend}: an alternative stress response of caring for others and seeking social support (more common in women). Chronic stress \textbf{suppresses the immune system} and raises the risk of heart disease (\textbf{psychophysiological illness}: real physical symptoms linked to stress).

\textbf{Coping}: \textbf{problem-focused} coping changes the stressor (make a study plan); \textbf{emotion-focused} coping manages the emotional reaction (talk to a friend, exercise, reframe). Perceived control, social support, exercise, and relaxation buffer stress.

\begin{q}{3}{5.1 GAS}{Prolonged stress depletes resources: exhaustion stage, illness risk rises.}
After months of caring for a sick parent while working two jobs, Elena keeps catching colds and feels drained. According to Selye, she is most likely in the
\chh{alarm stage}{resistance stage}{exhaustion stage}{tend-and-befriend stage}
\end{q}

\begin{q}{1}{5.1 Coping}{Directly addressing the stressor = problem-focused coping.}
Worried about failing chemistry, Raj meets with his teacher and makes a weekly study schedule. This is
\ch{problem-focused coping}{emotion-focused coping}{tend-and-befriend}{a psychophysiological illness}
\end{q}

\begin{q}{4}{5.1 Stressors}{Small, everyday irritations = daily hassles.}
A long commute, a broken phone screen, and a noisy roommate are best classified as
\chh{catastrophes}{adverse childhood experiences}{significant life changes}{daily hassles}
\end{q}

\section{Positive psychology (5.2)}

\textbf{Positive psychology} studies human flourishing: \textbf{subjective well-being} (life satisfaction and positive feelings), \textbf{positive emotions}, and \textbf{character strengths and virtues}. Six core virtues appear across cultures: wisdom, courage, humanity, justice, temperance, and transcendence. \textbf{Gratitude} practices (writing gratitude letters, counting blessings) reliably raise well-being. \textbf{Resilience} is the ability to recover from adversity; \textbf{posttraumatic growth} is positive change after struggling with trauma (deeper relationships, new priorities).

\begin{q}{2}{5.2 Posttraumatic growth}{Positive psychological change after trauma = posttraumatic growth.}
After surviving a serious accident, Omar says he has stronger friendships and a clearer sense of purpose than before. This illustrates
\chh{the exhaustion stage}{posttraumatic growth}{learned helplessness}{eustress}
\end{q}

\section{Explaining and classifying disorders (5.3)}

A \textbf{psychological disorder} is a syndrome marked by clinically significant disturbance in thought, emotion, or behavior that is typically \textbf{deviant} (atypical in context), \textbf{distressful}, and \textbf{dysfunctional} (interferes with daily life); danger to self or others can also be considered. Norms vary by culture and era.

The \textbf{DSM-5-TR} (American Psychiatric Association) lists diagnostic criteria; the \textbf{ICD} (World Health Organization) is used internationally. Diagnostic labels aid communication and research but can create \textbf{stigma} and bias how people are seen.

\begin{compare}{Perspective}{Cause of disorder}{Typical treatment}
Biological / medical & Genes, brain structure, neurotransmitters & Medication, brain stimulation \\
Psychodynamic & Unconscious conflicts from childhood & Psychoanalysis \\
Behavioral & Learned through conditioning and observation & Exposure, behavior modification \\
Cognitive & Maladaptive thoughts and beliefs & Cognitive therapy, CBT \\
Humanistic & Blocked growth, incongruent self-concept & Person-centered therapy \\
Sociocultural & Poverty, culture, discrimination, family & Community, family, group approaches \\
Evolutionary & Once-adaptive responses (e.g., fear) gone awry & --- \\
\end{compare}
The \textbf{biopsychosocial model} combines these. The \textbf{diathesis--stress model}: a predisposition (diathesis, often genetic) plus environmental stress triggers a disorder.

\begin{q}{2}{5.3 Diathesis-stress}{Genetic vulnerability + stressor triggers the disorder.}
A woman with a family history of depression develops it only after losing her job and her home. This best fits the
\ch{psychodynamic perspective}{diathesis--stress model}{humanistic perspective}{evolutionary perspective}
\end{q}

\begin{q}{4}{5.3 Perspectives}{Fear acquired by pairing with a bad experience = behavioral (learning) explanation.}
A psychologist explains a client's fear of dogs as the result of being bitten as a child and the fear being rewarded by escape. This is the
\chh{biological perspective}{psychodynamic perspective}{humanistic perspective}{behavioral perspective}
\end{q}

\section{Categories of disorders (5.4)}

\subsection{Neurodevelopmental}
\textbf{ADHD}: persistent inattention and/or hyperactivity-impulsivity beginning in childhood. \textbf{Autism spectrum disorder}: persistent deficits in social communication plus restricted, repetitive behaviors or interests, beginning in early childhood.

\subsection{Schizophrenia spectrum}
\begin{compare}{Symptom type}{Meaning}{Examples}
Positive & Presence of inappropriate behaviors & \textbf{Hallucinations} (false perceptions, often voices); \textbf{delusions} (false beliefs: persecution, grandeur); disorganized thinking and speech (word salad). \\
Negative & Absence of appropriate behaviors & \textbf{Flat affect} (no emotional expression), reduced speech, withdrawal, \textbf{catatonia} (motionless or odd postures). \\
\end{compare}
Explanations: the \textbf{dopamine hypothesis} (overactive dopamine activity; antipsychotics block dopamine receptors), enlarged ventricles and reduced cortical volume, genetic risk, prenatal factors, and diathesis--stress.

\subsection{Depressive and bipolar}
\textbf{Major depressive disorder}: at least two weeks of depressed mood or loss of interest/pleasure (\textbf{anhedonia}) plus symptoms such as sleep and appetite changes, fatigue, worthlessness. \textbf{Persistent depressive disorder}: milder but chronic depression lasting at least two years. Links: low serotonin and norepinephrine activity, a pessimistic explanatory style, rumination, and learned helplessness.
\textbf{Bipolar I}: at least one \textbf{manic} episode (elevated mood, little need for sleep, racing thoughts, grandiosity, risky behavior). \textbf{Bipolar II}: hypomanic (milder) episodes plus major depressive episodes.

\subsection{Anxiety}
\begin{terms}
Specific phobia & Intense, irrational fear of a specific object or situation (heights, spiders). \\
Agoraphobia & Fear of situations where escape might be hard (crowds, open spaces, leaving home). \\
Social anxiety disorder & Intense fear of being judged or embarrassed in social situations. \\
Panic disorder & Recurrent unexpected \textbf{panic attacks} (pounding heart, choking, sense of doom) and worry about future attacks. \\
Generalized anxiety disorder & Persistent, excessive, hard-to-control worry about many things. \\
\end{terms}

\subsection{Obsessive-compulsive, trauma, dissociative}
\textbf{OCD}: \textbf{obsessions} (intrusive, unwanted thoughts) and/or \textbf{compulsions} (repetitive behaviors, such as checking or washing, done to reduce anxiety). \textbf{Hoarding disorder} is a related condition. \textbf{PTSD}: after trauma, intrusive memories and nightmares, avoidance, negative changes in mood and thinking, and hyperarousal lasting more than a month. \textbf{Dissociative disorders} involve disrupted memory or identity: \textbf{dissociative amnesia} (inability to recall important personal information, sometimes with \textbf{fugue}, unplanned travel) and \textbf{dissociative identity disorder} (two or more distinct identities; diagnosis is controversial).

\subsection{Feeding and eating}
\textbf{Anorexia nervosa}: restriction of food intake leading to significantly low body weight, with intense fear of weight gain and distorted body image. \textbf{Bulimia nervosa}: recurrent binge eating followed by compensatory behaviors. \textbf{Binge-eating disorder}: recurrent binges without compensatory behaviors.

\subsection{Personality disorders}
Enduring, inflexible patterns that impair functioning, grouped in three clusters.
\begin{compare}{Cluster}{Theme}{Disorders}
A & Odd or eccentric & Paranoid (distrust), schizoid (detachment), schizotypal (odd beliefs and behavior). \\
B & Dramatic, emotional, erratic & \textbf{Antisocial} (disregard for others' rights, lack of remorse), \textbf{borderline} (unstable relationships, self-image, and mood; fear of abandonment), histrionic (attention seeking), narcissistic (grandiosity, need for admiration). \\
C & Anxious or fearful & Avoidant (fear of rejection), dependent (excessive need to be cared for), obsessive-compulsive personality (perfectionism, control). \\
\end{compare}

\begin{trap}
OCD vs.\ obsessive-compulsive \emph{personality} disorder: OCD involves unwanted intrusive thoughts and anxiety-reducing rituals the person usually recognizes as excessive; OCPD is a pervasive, rigid preoccupation with order and perfection that the person often sees as correct.
\end{trap}

\begin{q}{1}{5.4 Schizophrenia}{Hallucinations and delusions = positive symptoms.}
A man hears voices commenting on his actions and believes the government implanted a chip in his teeth. These are
\ch{positive symptoms of schizophrenia}{negative symptoms of schizophrenia}{symptoms of dissociative identity disorder}{symptoms of bipolar II}
\end{q}

\begin{q}{3}{5.4 Bipolar}{Days of little sleep, racing ideas, grandiosity, risky spending = mania.}
For a week, Nina sleeps two hours a night, talks rapidly about plans to start five companies, and spends her savings. These symptoms indicate a
\ch{panic attack}{major depressive episode}{manic episode}{dissociative fugue}
\end{q}

\begin{q}{2}{5.4 Anxiety}{Fear of places where escape is hard; staying home = agoraphobia.}
Lin has stopped going to malls, buses, and concerts because she fears she would be unable to escape if she panicked. She most likely has
\ch{specific phobia}{agoraphobia}{generalized anxiety disorder}{social anxiety disorder}
\end{q}

\begin{q}{4}{5.4 OCD}{Intrusive thoughts + anxiety-reducing rituals = OCD.}
Jack has intrusive thoughts that his house will burn down and checks the stove 30 times before leaving, which briefly relieves his anxiety. He most likely has
\ch{obsessive-compulsive personality disorder}{generalized anxiety disorder}{PTSD}{obsessive-compulsive disorder}
\end{q}

\begin{q}{1}{5.4 Personality}{Disregard for others' rights and no remorse = antisocial personality disorder.}
A man repeatedly cons elderly people out of their savings and shows no remorse. Which diagnosis is most consistent?
\ch{Antisocial personality disorder}{Avoidant personality disorder}{Schizoid personality disorder}{Dependent personality disorder}
\end{q}

\begin{q}{3}{5.4 Depression}{At least 2 years of milder chronic depression = persistent depressive disorder.}
For three years, Kayla has felt mildly down most days, with low energy and poor self-esteem, but never severely depressed. The best fit is
\ch{bipolar II}{major depressive disorder}{persistent depressive disorder}{generalized anxiety disorder}
\end{q}

\section{Treatment (5.5)}

\subsection{Psychological therapies}
\begin{compare}{Approach}{Goal}{Techniques}
Psychodynamic & Bring unconscious conflicts to awareness. & \textbf{Free association}, \textbf{dream analysis}, \textbf{interpretation} of \textbf{resistance} (blocking anxiety-laden material) and \textbf{transference} (feelings for others redirected onto the therapist). \\
Humanistic (person-centered) & Promote growth and self-acceptance. & Rogers's \textbf{nondirective} therapy: \textbf{active listening}, \textbf{unconditional positive regard}, genuineness, empathy. \\
Behavioral & Change maladaptive behavior directly. & \textbf{Exposure therapies}: \textbf{systematic desensitization} (relaxation + graded hierarchy), \textbf{flooding}, virtual-reality exposure. \textbf{Aversive conditioning} (pair unwanted behavior with unpleasantness). \textbf{Token economies} and behavior modification. \\
Cognitive & Change maladaptive thinking. & \textbf{Ellis's REBT}: dispute irrational ``musts.'' \textbf{Beck's cognitive therapy}: challenge automatic negative thoughts (the cognitive triad: self, world, future). \\
Cognitive-behavioral (CBT) & Change thoughts \emph{and} behaviors. & Thought records, homework, exposure. \textbf{DBT} (dialectical behavior therapy) adds emotion regulation and mindfulness; developed for borderline personality disorder. \\
\end{compare}
\textbf{Group therapy} and \textbf{family therapy} treat people within relationships and offer support. Many therapists use an \textbf{eclectic approach}. Effective therapy depends on a strong \textbf{therapeutic alliance}, evidence-based practice, and cultural sensitivity. Meta-analyses find therapy is generally effective, with some approaches best for specific disorders (exposure for phobias and OCD; CBT for depression and anxiety).

\subsection{Biomedical therapies}
\begin{terms}
Antidepressants & SSRIs (e.g., fluoxetine) increase serotonin availability by blocking reuptake; used for depression and anxiety disorders. \\
Antipsychotics & Block dopamine receptors; reduce positive symptoms. Long-term use of older drugs can cause \textbf{tardive dyskinesia} (involuntary facial and limb movements). \\
Anti-anxiety drugs & Benzodiazepines enhance GABA activity; risk of dependence. \\
Mood stabilizers & \textbf{Lithium} for bipolar disorder. \\
Stimulants & Used to treat ADHD (increase dopamine and norepinephrine activity). \\
ECT & Electroconvulsive therapy: brief electric current under anesthesia; effective for severe, treatment-resistant depression. \\
TMS & Transcranial magnetic stimulation: magnetic pulses to the cortex; for depression. \\
Psychosurgery & Removing or destroying brain tissue (historical lobotomy); now rare and highly targeted. \\
\end{terms}

\begin{q}{2}{5.5 Psychodynamic}{Redirecting feelings about a parent onto the therapist = transference.}
A client begins reacting to her therapist with the same anger she feels toward her controlling father. A psychodynamic therapist would call this
\chh{resistance}{transference}{free association}{counterconditioning}
\end{q}

\begin{q}{4}{5.5 Behavioral}{Relaxation paired with a graded fear hierarchy = systematic desensitization.}
A therapist teaches a client deep relaxation, then has her imagine a series of increasingly frightening flying scenarios while staying relaxed. This is
\chh{flooding}{aversive conditioning}{a token economy}{systematic desensitization}
\end{q}

\begin{q}{1}{5.5 Cognitive}{Challenging automatic negative thoughts about self, world, and future = Beck's cognitive therapy.}
A therapist helps a depressed client notice the automatic thought ``I'm a failure'' and test it against evidence. This approach is most associated with
\ch{Beck's cognitive therapy}{Rogers's person-centered therapy}{Freud's psychoanalysis}{flooding}
\end{q}

\begin{q}{3}{5.5 Biomedical}{Lithium is the classic mood stabilizer for bipolar disorder.}
Which treatment is most closely associated with stabilizing mood in bipolar disorder?
\chh{SSRIs}{Antipsychotics}{Lithium}{Benzodiazepines}
\end{q}

\begin{q}{2}{5.5 Humanistic}{Nondirective, reflective listening with acceptance = person-centered (Rogers).}
A therapist mainly restates the client's feelings (``It sounds like you feel overlooked'') and avoids giving advice. This is
\ch{psychodynamic interpretation}{person-centered therapy using active listening}{rational emotive behavior therapy}{aversive conditioning}
\end{q}

\makekey
\end{document}
